\fontsize{12pt}{14pt} \selectfont
\begin{doublespace}
\chapter{Conclusion}Containers have become a cornerstone of modern software deployment, valued for their efficiency, portability, and scalability. However, ensuring that containers are securely confined remains a complex and evolving challenge. This thesis introduced Houdini, a modular framework designed to empirically evaluate container security confinement across a wide range of configurations. By executing controlled security tests, referred to as tricks, within a QEMU-based virtual machine, Houdini enables direct observation of whether isolation mechanisms such as namespaces, control groups, seccomp, AppArmor, and Linux capabilities are enforced as expected during runtime.

The case study chapter demonstrated Houdini's ability to systematically uncover gaps between intended and actual security guarantees in containerized environments. Through a series of examples, we showed how Houdini can identify vulnerabilities that persist across multiple versions of runtime components, as illustrated by the CVE-2024-21626 example. We also analyzed the impact of privilege and namespace flags on host interaction and detected configuration drift resulting from changes in runtime behavior or policy enforcement. These results highlight the value of empirical and reproducible testing for improving container confinement.

Significantly, Houdini’s applications extend beyond the specific scenarios examined in this thesis. Its modular and extensible design allows users to define and share new tricks for testing a wide range of confinement violations, misconfigurations, or security behaviors. As a general-purpose framework, Houdini is well suited to support both researchers and practitioners in addressing emerging security challenges in containerized systems.
Houdini’s design also supports continuous security testing in modern software development workflows. Because it is both safe and scriptable, Houdini can be integrated into automated testing systems, such as continuous integration and deployment pipelines, to provide ongoing validation of container security. Whenever any part of the container stack changes, such as the kernel, Docker engine, or a security policy, Houdini can rerun a suite of empirical tests to ensure that the confinement guarantees remain intact.

In this context, Houdini functions as a form of security regression testing. By systematically repeating empirical tests across different software versions and configurations, Houdini helps detect enforcement failures and configuration drift that might otherwise go unnoticed. This ensures that the security boundaries surrounding containers remain strong as platforms evolve.

To encourage further research and community adoption, the source code for Houdini has been made publicly available\footnote{\url{https://github.com/CU-Security-Labs/houdini}}, along with a set of initial test cases. We hope this enables others to expand Houdini’s capabilities, contribute new security tests, and collectively improve the state of container security evaluation.

\end{doublespace}